\chapter*{General Introduction}
\addcontentsline{toc}{chapter}{General Introduction}

Industrial performance is not improved by isolated algorithms; it is improved when reliable information reaches the right decision-maker at the right time. In a continuous production setting, this means linking process understanding, quality management, information flow and operator action. This internship is consequently positioned as an \textbf{Industrial Engineering digitalization and supervision project using data and AI technologies}, rather than as an artificial-intelligence exercise detached from production.

At OCP Jorf Lasfar, the soluble monoammonium phosphate production chain comprises successive chemical and physical transformations, from reactant preparation and neutralization to crystallization, solid--liquid separation, drying and conditioning. The project narrows this chain to the drying stage because final moisture is strongly influenced there, moisture is a direct quality indicator, and abnormal drying behaviour can rapidly affect conformity. The industrial difficulty is temporal: key process measurements are continuously available to supervision, while laboratory quality results become available only about once every two hours. Operators therefore face a period in which the most recent laboratory result may no longer fully represent the current state.

The project examines whether this visibility gap can be reduced through a digital decision-support chain. A moisture soft sensor estimates final moisture from process behaviour; a novelty detector identifies operating observations that depart from a reference regime; and a diagnostic layer translates deviations into contributors, subsystems, cautious probable diagnoses and verification actions. A Python service replays prototype rows on a deliberately selected five-second cycle, PostgreSQL centralizes the data, and a Power BI DirectQuery report provides an operator-oriented interface. The five-second interval is a practical simulation and visualization frequency, not the native measurement frequency of the plant instrumentation.

The scope boundary is fundamental. The industrial process, observed PCS7 screens, laboratory workflow and selected variables are real anchors. The final development and demonstration datasets, however, are synthetic prototypes designed to reproduce realistic multi-rate behaviour. There is no direct connection to the live PCS7 system, no autonomous actuation, no closed-loop production control and no claim of industrial-scale validation. A five-second replay demonstrates the software architecture; it does not certify plant performance.

The adopted workflow follows the logic of an industrial improvement project: understand the process and the decision latency; establish a data contract; construct and audit the dataset; prevent temporal leakage; compare models; integrate inference with a traceable database; design a supervision interface; and validate the complete chain. Python, scikit-learn, PostgreSQL and Power BI are enabling tools within that workflow.

Three questions guide the study. First, can fast process history provide a useful estimate of a quality variable whose direct observation is sparse? Second, can abnormal operation be communicated in a form that supports physical verification rather than presenting an opaque model score? Third, can the complete information chain operate frequently enough for near-real-time supervision while preserving traceability and failure visibility? These questions connect modelling performance with the wider production system.

The project also adopts explicit information categories. A \emph{measurement} is a value acquired from the process or laboratory source. A \emph{prediction} is a model output derived from a defined historical window. An \emph{anomaly} is a statistical departure from the learned reference, not necessarily a fault. A \emph{diagnosis} is a ranked, rule-supported interpretation, not proof of cause. A \emph{recommendation} is a verification prompt, not an instruction to change a setpoint. Keeping these categories separate reduces automation bias and provides a basis for future operator training.

The expected operational benefit is consequently expressed in terms of visibility and response preparation, not as an unverified economic saving. A continuous estimate may reveal a developing quality tendency between lab results; an anomaly event may direct attention to the most relevant subsystem; and centralized history may shorten the time needed to reconstruct what changed. Quantified benefits such as reduced off-specification production, lower energy consumption or shorter response time would require baseline plant data and a controlled industrial evaluation that are outside the initiation-internship scope.

Finally, the report treats reproducibility as part of engineering quality. Dataset cadence and sample counts are audited, metrics are read from saved training records, model and feature-schema versions are persisted, database writes are idempotent, validation commands are recorded and the report figures are tied to repository evidence. This discipline makes the current limitations inspectable and prepares a better hand-off for a later industrial study.

This report is organized in eight chapters. Chapters 1 and 2 establish the concise industrial context and the soluble MAP process. Chapter 3 formalizes the problem, scope and objectives. Chapter 4 documents data construction, preparation and feature engineering. Chapter 5 presents the predictive and diagnostic models. Chapters 6 and 7 describe the prototype replay architecture and the Power BI dashboard. Chapter 8 consolidates results, tests, limitations and the future industrialization roadmap.

\begin{limitbox}[Interpretation of every result]
All numerical model and runtime results in this report concern the repository's prototype/synthetic datasets. They demonstrate technical feasibility under controlled replay conditions, not production-grade accuracy, reliability, safety integrity or economic benefit.
\end{limitbox}
